\documentclass[11pt,a4paper]{article}

\usepackage[margin=25mm]{geometry}
\usepackage[T1]{fontenc}
\usepackage{lmodern}
\usepackage{amsmath,amssymb,amsthm}

\theoremstyle{definition}
\newtheorem{definition}{Definition}
\newtheorem{problem}{Problem}

\setlength{\parindent}{0pt}
\setlength{\parskip}{6pt}

\begin{document}

\begin{center}
  {\Large\bfseries Davenport--Schinzel Sequences: The Pigeonhole Threshold}
\end{center}

\section{Definition}
\begin{definition}[Order-$s$ Davenport--Schinzel sequence]
Let $n,s\geq 1$ and $[n]=\{1,\ldots,n\}$. A sequence
$U=(u_1,\ldots,u_m)$ over $[n]$ is an order-$s$ Davenport--Schinzel sequence if
\begin{enumerate}
  \item $u_i\neq u_{i+1}$ for every $i$; and
  \item for every two distinct symbols $a,b$, the sequence $U$ contains no
  alternating subsequence
  \[
    a,b,a,b,\ldots
  \]
  of length $s+2$.
\end{enumerate}
We write $U\in\mathrm{DS}(n,s)$.
\end{definition}

\begin{definition}[Extremal function]
\begin{equation}
  \lambda_s(n)
  =\max\bigl\{|U|:U\in\mathrm{DS}(n,s)\bigr\}.
\end{equation}
\end{definition}

For all $n$ and $s$, the pigeonhole argument gives
\begin{equation}
  \lambda_s(n)\leq\binom{n}{2}(s+1).
\end{equation}

\section{Current best result}
This bound is asymptotically tight when
\begin{equation}
  \frac{s}{\sqrt n}\longrightarrow\infty.
\end{equation}
In this regime,
\begin{equation}
  \lambda_s(n)\sim\binom{n}{2}s.
\end{equation}

\section{Problem Formulation}

Determine the smallest growth rate of $s=s(n)$ for which
\begin{equation}
  \lambda_s(n)\sim\binom{n}{2}s.
\end{equation}
In particular, determine whether this equivalence can hold for some
$s=o(\sqrt n)$.
\end{document}